% Deliberately broken example manuscript for testing the
% manuscript-preflight skill and check_manuscript.py.
% The topic, authors, and references are entirely fictional.
% Each problem below is annotated with the check that should catch it.
\documentclass[11pt]{article}

\usepackage{graphicx}
\usepackage{epsfig}      % deprecated package
\usepackage{subfigure}   % deprecated package
\usepackage{amsmath}

% TODO: check with coauthors whether to keep the stability section  <- private note in comment

\title{Convergence of an Adaptive Mesh Refinement Scheme for a Model Diffusion Problem}
\author{A. Author\\Department of Applied Mathematics, Example University}
\date{}

\begin{document}

\maketitle

\begin{abstract}
We analyze the convergence of an AMR scheme for a model diffusion problem.
% "AMR" used before being defined anywhere -> abbreviation audit
We prove that the scheme achieves a convergence rate of 2.1 in the energy norm.
% Abstract says 2.1, body says 1.8 -> abstract/body consistency
\end{abstract}

\section{Introduction}
\label{sec1} % opaque label

Adaptive discretization methods have been studied extensively
\cite{doe2015}. Recent work \cite{smith2018,nonexistent2024} extends these
ideas to time-dependent problems.
% nonexistent2024 is not in refs.bib -> undefined citation

As shown in Fig. 2, the error decreases monotonically with refinement level.
% hardcoded figure number instead of \ref -> semantic check

{\bf Definition 1.} A mesh is \emph{admissible} if every element satisfies
the minimum-angle condition and neighboring elements differ by at most one
refinement level.
% bold-text pseudo-environment -> should be a real definition environment

\section{Model Problem}
\label{sec:model}

The model equation reads
$$
\partial_t u = \nabla \cdot \left( \kappa \nabla u \right) + f,
$$
% $$...$$ plain-TeX display math, and the equation is unnumbered/unlabeled
where the meaning of $\kappa$ and $f$ is not introduced anywhere nearby.
% symbol-introduction audit

\begin{equation}
\label{sec:model} % duplicate label (already used for the section above)
\eta_K^2 = h_{\rm eff}^2 \| f + \nabla \cdot ( \kappa \nabla u_h ) \|_K^2
% math-mode {\rm ...} -> INFO with \mathrm suggestion (not a WARN)
\end{equation}

The estimator is reliable and efficient, see Eq.~\eqref{eq:efficiency}.
% undefined reference: no \label{eq:efficiency} exists

\section{Results}
\label{sec:results}

The scheme achieves a convergence rate of 1.8 in the energy norm.

\begin{table}[t]
\caption{Parameters used in the numerical experiments.}
\label{tab:parameters}
\includegraphics[width=\columnwidth]{/home/username/plots/table_screenshot.png}
% table shipped as a screenshot AND via an absolute local path
\end{table}

\begin{figure}[t]
\includegraphics[width=\columnwidth]{fig2_final_v3.pdf}
% on disk the file is Fig2_final_v3.pdf -> case-mismatch ERROR (arXiv-fatal);
% also an opaque, versioned figure filename
\caption{Results.} % caption cannot stand alone
\label{fig2} % opaque label
\end{figure}

Figure~\ref{fig2} shows the main result.

\section{Conclusion}
\label{sec:conclusion} % this label is never referenced -> unused-label INFO

We conclude that the adaptive scheme converges at the optimal rate for the
model problem.
\todo{Tighten the outlook paragraph before submission.}
% draft note command left in the source -> draft-artifact check

\bibliographystyle{plain}
\bibliography{refs}

\end{document}
